\documentclass[11pt,a4paper]{article}
\usepackage[italian]{babel}
\usepackage[margin=2.5cm]{geometry}
\usepackage{parskip}
\usepackage{amsmath, amssymb, amsthm}
\usepackage{booktabs}
\usepackage{array}
\usepackage{xcolor}
\usepackage{needspace}
\usepackage{enumitem}
\usepackage{listings}
\usepackage{tikz}
\usepackage[most]{tcolorbox}
\usepackage[hidelinks]{hyperref}
\usetikzlibrary{decorations.pathreplacing, shapes.geometric, arrows.meta, positioning}

% Niente vedove/orfane: i paragrafi non lasciano righe isolate a fine/inizio pagina
\widowpenalty=10000
\clubpenalty=10000

% Elenchi più compatti
\setlist{itemsep=2pt, parsep=0pt, topsep=4pt, partopsep=0pt}

% --- Riquadri con bordo nero, senza numero (si spezzano tra due pagine) ---
% Uso: \begin{definizione} ... \end{definizione}
%      \begin{teorema}[Nome] ... \end{teorema}  (il nome è facoltativo)
\tcbset{nota/.style={enhanced jigsaw, breakable, colback=white,
  colframe=black, boxrule=0.5pt, sharp corners}}
\NewTColorBox{definizione}{}{nota,
  before upper={\textbf{Definizione.}\ }}
\NewTColorBox{teorema}{o}{nota,
  before upper={\textbf{Teorema\IfValueT{#1}{ (#1)}.}\ }}
\NewTColorBox{proposizione}{o}{nota,
  before upper={\textbf{Proposizione\IfValueT{#1}{ (#1)}.}\ }}
\NewTColorBox{proprieta}{o}{nota,
  before upper={\textbf{Proprietà\IfValueT{#1}{ (#1)}.}\ }}
\NewTColorBox{assioma}{o}{nota, boxrule=1pt,
  before upper={\textbf{Assioma\IfValueT{#1}{ (#1)}.}\ }}

% --- Ambienti semplici (senza riquadro) ---
% Il titolo sta in cima al blocco, anche se il contenuto inizia con un riquadro o
% con codice; e non resta mai da solo in fondo a una pagina.
% Uso: \begin{esempio}[Titolo facoltativo] ... \end{esempio}
\newcommand{\newrunin}[2]{%
  \NewDocumentEnvironment{#1}{o}{%
    \par\Needspace{4\baselineskip}\addvspace{\medskipamount}\noindent
    \textbf{#2}\IfValueT{##1}{ (##1)}\textbf{.}\ \ignorespaces}%
  {\par\addvspace{\medskipamount}}}
\newrunin{esempio}{Esempio}
\newrunin{esercizio}{Esercizio}
\newrunin{osservazione}{Osservazione}

% V e F nelle tavole di verità
\newcommand{\V}{\textsf{V}}
\newcommand{\F}{\textsf{F}}

% --- Codice C ---
% Uso: \begin{codice} ... \end{codice}      codice C numerato
%      \begin{comando} ... \end{comando}    comandi da terminale
%      \begin{risultato} ... \end{risultato} output di un programma
%      \lstinline|printf("ciao");|           codice dentro al testo
\lstdefinestyle{c}{
  language=C,
  basicstyle=\ttfamily\small,
  keywordstyle=\color{blue}\bfseries,
  commentstyle=\color{gray}\itshape,
  stringstyle=\color{red!70!black},
  numbers=left,
  numberstyle=\tiny\color{gray},
  numbersep=8pt,
  columns=flexible,
  keepspaces=true,
  breaklines=true,
  showstringspaces=false,
  tabsize=4,
  morekeywords={include, define},
  % lettere accentate nei commenti
  literate={à}{{\`a}}1 {è}{{\`e}}1 {é}{{\'e}}1 {ì}{{\`i}}1
           {ò}{{\`o}}1 {ù}{{\`u}}1 {È}{{\`E}}1 {À}{{\`A}}1
}
\lstset{style=c}
\newtcblisting{codice}{listing only, nota, left=6mm,
  listing options={style=c}}
\newtcblisting{comando}{listing only, nota,
  listing options={basicstyle=\ttfamily\small, numbers=none, language={},
    columns=flexible, keepspaces=true, breaklines=true}}
\newtcblisting{risultato}{listing only, enhanced jigsaw, breakable,
  colback=black!5, colframe=black, boxrule=0.5pt, sharp corners,
  listing options={basicstyle=\ttfamily\small, numbers=none, language={},
    columns=flexible, keepspaces=true, breaklines=true}}

% --- Diagrammi di flusso (simboli standard) ---
\tikzset{
  terminale/.style = {ellipse, draw, minimum width=2.5cm, minimum height=0.9cm},
  io/.style        = {trapezium, trapezium left angle=70, trapezium right angle=110,
                      draw, minimum width=2.5cm, minimum height=0.9cm},
  processo/.style  = {rectangle, draw, minimum width=2.5cm, minimum height=0.9cm},
  decisione/.style = {diamond, draw, aspect=2, minimum width=2.5cm},
  freccia/.style   = {thick, -{Stealth}}
}

\title{Geometria analitica: la retta}
\author{Marco Cerbella}
\date{Lezione 3 -- 1 ottobre 2026}

\begin{document}
\maketitle

\section{Equazione cartesiana della retta}

Un'equazione algebrica di primo grado in due variabili
\[
ax + by + c = 0, \qquad a, b, c \in \mathbb{R} \tag{1}
\]
rappresenta una retta in posizione generica nel piano: \emph{equazione cartesiana generale della retta in forma implicita}. Fissati i coefficienti $a, b, c$ è fissata la retta, nel senso che tutti i punti della retta hanno coordinate $(x, y)$ che soddisfano l'equazione.

\begin{esempio}
I punti della retta $y - x = 0$ devono avere ascissa uguale all'ordinata. $P = (1, 2)$ la soddisfa? No ($2 - 1 \neq 0$). $P = (5, 5)$ la soddisfa? Sì.
\end{esempio}

\section{Forma esplicita}

Se $b \neq 0$ si isola il termine in $y$ e si divide per $b$:
\[
ax + by + c = 0 \implies by = -ax - c \implies y = -\frac{a}{b}x - \frac{c}{b}.
\]
Si ottiene l'equazione in \emph{forma esplicita}
\[
y = mx + q, \qquad m, q \in \mathbb{R}, \tag{2}
\]
con $m = -\frac{a}{b}$ e $q = -\frac{c}{b}$ (queste relazioni permettono di passare da una forma all'altra).

L'equazione (1) può rappresentare una qualunque retta del piano. L'equazione (2) rappresenta tutte le rette \textbf{tranne} quelle parallele all'asse delle $y$.

L'equazione (2) è il grafico di una funzione: $f : D \subseteq \mathbb{R} \to \mathbb{R}$, $x \mapsto f(x) = mx + q$ con $D \subseteq \mathbb{R}$ e
\[
\mathrm{Gr}(f) = \{(x, f(x)) \mid x \in D\} = \{(x, mx + q) \mid x \in D\}.
\]

\section{Rette parallele agli assi}

\begin{itemize}
    \item Una retta parallela all'asse delle ordinate ha equazione $x = x_0$, perché è costituita da infiniti punti aventi la stessa ascissa, cioè i punti del tipo $(x_0, y)$, $y \in \mathbb{R}$, con $x_0$ costante. \textbf{Non può essere il grafico di una funzione}: allo stesso $x_0$ corrispondono infiniti valori di $y$.
    \item Una retta parallela all'asse delle ascisse ha equazione $y = y_0$, perché è costituita da infiniti punti aventi la stessa ordinata, cioè i punti del tipo $(x, y_0)$, $x \in \mathbb{R}$, con $y_0$ costante. È il grafico della funzione costante $f : \mathbb{R} \to \mathbb{R}$, $x \mapsto y_0$.
\end{itemize}

\section{Pendenza e intercetta}

Per la retta $y = mx + q$:
\begin{itemize}
    \item $m$ è la \emph{pendenza}: $m = \dfrac{y_2 - y_1}{x_2 - x_1}$, con $(x_1, y_1), (x_2, y_2)$ punti della retta;
    \item $m = \tan\theta$, dove $\theta$ è l'angolo formato tra l'asse delle ascisse e la retta;
    \item $m = 0$: la retta è parallela all'asse delle ascisse ($y = q$);
    \item $q$ è l'intersezione della retta con l'asse delle ordinate;
    \item $q = 0$: la retta passa per l'origine.
\end{itemize}

\section{Zeri e intersezione con l'asse $x$}

Le soluzioni di $ax + b = 0$, con $a, b \in \mathbb{R}$, $a \neq 0$, sono graficamente le soluzioni del sistema
\[
\begin{cases} y = 0 \\ y = ax + b \end{cases}
\implies \text{soluzione } P = \left(-\frac{b}{a}, 0\right) = (x_0, 0).
\]
Cioè, data $f : \mathbb{R} \to \mathbb{R}$, $f(x) = ax + b$, la domanda ``per quali $x \in \mathbb{R}$ si ha $f(x) = 0$?'' ha come risposta l'ascissa del punto in cui il grafico interseca l'asse $x$.

\end{document}
